\documentclass[11pt]{article}
\usepackage[margin=0.75in]{geometry}
\usepackage{amsmath,amssymb,booktabs}
\usepackage[hidelinks]{hyperref}
\setlength{\parindent}{0pt}
\setlength{\parskip}{5pt}
\title{\vspace{-1.5em}Opening phase as a boundary interaction\\
\large A verified \(46\,173\)-letter superpermutation on eight symbols}
\author{}
\date{20 October 2026}
\begin{document}
\maketitle
\vspace{-2.5em}

The improvement retains all sixteen closed trails of Jakub Halfar's
selected instance of Jay Pantone's construction: four large trails arising
from four \(27\)-loop walks, and twelve trails completing the remaining
source loops. No trail is split or deleted, and no opening drops a
permutation occurrence. Only trail order and opening positions change.
The input selection and the search program are Halfar's; the experiment
reported here is a boundary reassembly, not a new all-order construction.

\section*{1. Phase determines two literal boundaries}
For the general spelling identity, let \(n\geq4\), put \(h=n-3\), and
let trail \(i\) have cyclic spelling length \(R_i\).
An admissible opening \(\alpha\) specifies a cyclic start position and a
gap \(g_{i,\alpha}\) between the adjacent permutation windows at that cut.
The opened piece has length
\[
R_i+n-g_{i,\alpha}
 =R_i+h+a_i(\alpha),
\qquad
a_i(\alpha)=3-g_{i,\alpha}.
\]
It exposes a starting \(h\)-word \(S_{i,\alpha}\) and an ending
\(h\)-word \(E_{i,\alpha}\).
\emph{Phase} here means the full cyclic position, not just a label modulo
\(h\): equal coarse clock phases need not give equal boundary words.
For this experiment, weight-\(3\) openings have \(a=0\), while
weight-\(2\) openings have \(a=1\).

\section*{2. One common matrix; different boundary states}
Let \(\mathcal B_n\) be the ordered \(h\)-tuples of distinct alphabet
symbols. Its size is \(n!/3!\), hence \(6\,720\) at \(n=8\).
For \(E,S\in\mathcal B_n\), define
\[
\begin{gathered}
\operatorname{ov}(E,S)
 =\max\{k\in\{0,\ldots,h\}:\operatorname{suffix}_k(E)
                         =\operatorname{prefix}_k(S)\},\\
D[E,S]=h-\operatorname{ov}(E,S).
\end{gathered}
\]
The length-zero prefix and suffix are both the empty word.
Thus \(D\) is a directed boundary-cost matrix, common to every pair of
trails. It need not be stored densely. If \(\mathbf e_E\) denotes the
one-hot vector for boundary word \(E\), the exact cost of a join is
\[
\mathbf e_{E_{i,\alpha}}^{T}D\,
\mathbf e_{S_{j,\beta}}.
\]
Collect the endpoint vectors as columns of matrices
\(\mathsf E_i\) and \(\mathsf S_i\). The interaction between opening
choices of two trails is then
\[
\boxed{C_{ij}=\mathsf E_i^{T}D\,\mathsf S_j},
\qquad
C_{ij}=
\begin{pmatrix}
D[E_{i,1},S_{j,1}] & D[E_{i,1},S_{j,2}] & \cdots\\
D[E_{i,2},S_{j,1}] & D[E_{i,2},S_{j,2}] & \cdots\\
\vdots & \vdots & \ddots
\end{pmatrix}.
\]
For order \(\sigma\) on the sixteen trails, let \(x_i\) be a one-hot
opening-choice vector, and let \(a_i\) be the vector of opening costs.
The exact spelling identity for these unsplit, non-dropping pieces is
\[
\boxed{
L=h+\sum_{i=1}^{16}R_i
 +\sum_{i=1}^{16}a_i^{T}x_i
 +\sum_{k=1}^{15}
 x_{\sigma(k)}^{T}C_{\sigma(k),\sigma(k+1)}x_{\sigma(k+1)}.
}
\]
The two ends of a trail are coupled through the \emph{same} \(x_i\).
Optimizing its incoming and outgoing boundaries independently would
therefore admit choices that no single opening realizes.

\section*{3. An actual expanded matrix from the new word}
The three weight-\(2\) openings occur consecutively, in the order
\(14,15,13\). Trail numbers below are the zero-based IDs in the saved
plan; offsets are positions in their cyclic spellings.
\[
\begin{array}{c|r|c|c|c}
\text{trail}&\text{offset}&g&S&E\\ \hline
14&6538&2&21750&17503\\
15&8616&2&17503&75032\\
13&3415&2&03254&32541
\end{array}
\]
Using row order \(E_{14},E_{15},E_{13}\) and column order
\(S_{14},S_{15},S_{13}\), the literal overlap and cost matrices are
\[
O=
\begin{pmatrix}
0&5&2\\
1&0&3\\
0&1&0
\end{pmatrix},
\qquad
M=5\mathbf1\mathbf1^{T}-O
=
\begin{pmatrix}
5&\boxed{0}&3\\
4&5&\boxed{2}\\
5&4&5
\end{pmatrix}.
\]
Here \(\mathbf1=(1,1,1)^T\). These are exact boundary costs for the
displayed openings, not hypothetical examples. Diagonal entries are
shown for completeness; revisiting a trail is not allowed.

The two joins actually used in this subchain are the matrix elements
\[
\begin{pmatrix}1&0&0\end{pmatrix}
M\begin{pmatrix}0\\1\\0\end{pmatrix}=0,
\qquad
\begin{pmatrix}0&1&0\end{pmatrix}
M\begin{pmatrix}0\\0\\1\end{pmatrix}=2.
\]
The zero is the literal equality \(E_{14}=S_{15}=17503\).
The next join shares the three-symbol boundary \(032\).
This small matrix illustrates the boundary mechanism; the net saving
must still be accounted for over the entire sixteen-trail assembly.

\section*{4. Three extra opening symbols buy four joining symbols}
Write \(A=\sum_i a_i^Tx_i\), and write \(J\) for the sum of the fifteen
join costs. The fixed cyclic spelling length is
\(\sum_iR_i=46\,144=46\,080+64\), so \(h+\sum_iR_i=46\,149\).
The exact comparison is
\[
\begin{array}{l|r|r}
&\text{Halfar reference}&\text{new assembly}\\ \hline
h+\sum_iR_i&46\,149&46\,149\\
A&0&3\\
J&25&21\\ \hline
L&46\,174&46\,173
\end{array}
\]
For the new plan the fifteen joins, in order, expand to
\[
J=1+1+3+1+1+2+2+0+2+1+1+1+3+1+1=21.
\]
Consequently
\[
\boxed{\Delta L=\Delta A+\Delta J=3+(21-25)=-1.}
\]
The gain is collective: the accounting does not assign an independent
saving to each of the three changed openings.

\section*{5. Occupation and scope}
Every required permutation remains covered by the complete trail
pieces, and overlapping equal boundary strings preserves their
substrings. This preserves final coverage, not the intermediate
Bellman states or the count of repeated occurrences.
Two independent literal scanners accepted length \(46\,173\) and all
\(8!=40\,320\) permutations. The new word has \(56\) repeated permutation
occurrences, versus \(57\) in the reference, and both have \(5\,790\)
dirty windows.

In our Albion picture, this is an occupation-preserving change of
boundary phase. The matrix formulation explains the coupled cost;
Halfar's existing randomized \texttt{trailsearch} found the discrete
choices. Fourier diagonalization of an internal clock does not
automatically diagonalize \(D\), and no general analytic selection rule
or minimality theorem is claimed.

{\small\textit{Reproducibility.}
The adjacent \texttt{../PROBE27.md} records commands, attribution and
verification. The plan is \texttt{../n8-46173.plan}; the word is
\texttt{../n8-46173.txt}.}
\end{document}
